% संपूर्ण मराठी ग्रंथातील प्रकरण 40: त्रिमूल्य तर्कशास्त्रे
% हा संपादनयोग्य स्रोतखंड आहे. संकलनासाठी संपूर्ण स्रोतसंग्रहातील
% अक्षररूपे, पूर्वभाग, चिन्हव्याख्या आणि आकृती-स्रोत आवश्यक आहेत.
% मूळ: Open Logic Project; मजकूर आणि रूपांतर CC BY 4.0.
% यंत्रानुवाद, दुरुस्ती आणि तपासणी: OpenAI Codex —
% GPT-5.6 Sol आणि GPT-6 Sol, दोन्ही Ultra effort.
% दुरुस्ती-पुनरावलोकन: GPT-6.1 Sol, Ultra effort.
\chapter{त्रिमूल्य तर्कशास्त्रे}\label{mvl:thr:chap}









\section{परिचय}\label{mvl:thr:int:sec}

$\True$ आणि $\False$ या मूल्यांच्या जोडीला आणखी एक मूल्य~$\Undef$
जोडले, तर त्रिमूल्य तर्कशास्त्र मिळते. हे एकच अतिरिक्त
सत्यतामूल्य असले, तरी $\lnot$, $\land$, $\lor$ आणि $\lif$
यांसाठी सत्यता-फलने परिभाषित करण्याचे अनेक पर्याय असतात.
मग एखादे तर्कशास्त्र या सत्यता-फलनांचा कोणताही संयोग वापरू
शकते; तसेच फक्त $\True$ ला निर्दिष्ट मानायचे की $\True$
आणि~$\Undef$ या दोन्हींना, हाही पर्याय असतो.

येथे आपण सर्वाधिक परिचित त्रिमूल्य तर्कशास्त्रांपैकी काही,
त्यांमागील प्रेरणा आणि त्यांचे काही गुणधर्म मांडतो.













\section{\L ukasiewicz यांचे तर्कशास्त्र}\label{mvl:thr:luk:sec}

बहुमूल्य तर्कशास्त्राचे प्रकाशित झालेले आणि सविस्तर मांडलेले
सर्वांत आरंभीच्या प्रस्तावांपैकी एक पोलिश तत्त्वज्ञ Jan \L ukasiewicz
यांनी 1920 मध्ये मांडला.\footnote{Jan \L ukasiewicz,
``O logice trójwartościowej, \emph{Ruch Filozoficzny} 5
(1920), पृ. 170--171; 19 जून 1920 च्या व्याख्यानाचा वृत्तांत.
मूळ अंकाचा स्कॅन: \url{https://kpbc.umk.pl/Content/247342/PDF/EOD_OPEN_006_12_nr_09.pdf}.} अरिस्टॉटलच्या सागरी युद्धाच्या समस्येतून
\L ukasiewicz यांना प्रेरणा मिळाली: \emph{आज} ``उद्या सागरी युद्ध
होईल'' हे विधान सत्यही नाही आणि असत्यही नाही असे दिसते; त्याचे
सत्यतामूल्य अजून निश्चित झालेले नाही. अशा ``भविष्यातील
परिस्थितीवश'' विधानांसाठी \L ukasiewicz यांनी तिसरे सत्यतामूल्य
आणण्याचा प्रस्ताव मांडला.
\begin{quote}
पुढील वर्षी एखाद्या ठराविक क्षणी, उदाहरणार्थ 21 डिसेंबरला दुपारी बारा वाजता,
मी वॉर्सामध्ये असेन, हे सध्या होकारार्थी किंवा नकारार्थी अशा
कोणत्याही प्रकारे निश्चित झालेले नाही, असे मी विरोधाभासाशिवाय
गृहीत धरू शकतो. म्हणून त्या वेळी मी वॉर्सामध्ये असेन हे संभाव्य
आहे, पण अनिवार्य नाही. या गृहितकावर ``पुढील वर्षी 21 डिसेंबरला
दुपारी बारा वाजता मी वॉर्सामध्ये असेन'' हे विधान सध्या सत्यही नाही आणि
असत्यही नाही. कारण ते आताच सत्य असेल, तर भविष्यात मी
वॉर्सामध्ये असणे अनिवार्य ठरेल; हे गृहितकाशी विसंगत आहे.
दुसरीकडे ते आताच असत्य असेल, तर भविष्यात मी वॉर्सामध्ये असणे
अशक्य ठरेल; हेही गृहितकाशी विसंगत आहे. म्हणून विचाराधीन
विधान या क्षणी सत्यही नाही आणि असत्यही नाही; त्याला ``0''
म्हणजे असत्यता आणि ``1'' म्हणजे सत्यता यांपेक्षा वेगळे तिसरे
मूल्य असले पाहिजे. या मूल्याला आपण ``$\frac{1}{2}$.'' असे
दर्शवू शकतो. ते ``संभाव्य'' याचे प्रतिनिधित्व करते आणि
``सत्य'' व ``असत्य'' यांच्याबरोबर तिसरे मूल्य म्हणून सामील होते.
\end{quote}
\L ukasiewicz यांच्या तिसऱ्या सत्यतामूल्यासाठी आपण $\Undef$
वापरू.\footnote{\L ukasiewicz येथे ``संभाव्य'' हा शब्द आज रूढ
नसलेल्या अर्थाने, म्हणजे संभाव्य पण अनिवार्य नाही, या अर्थाने
वापरतात.}

या अर्थानुसार $\lnot$, $\land$ आणि $\lor$ या संयोजकांची
सत्यता-फलने ठरवणे सोपे आहे: भविष्यातील परिस्थितीवश विधानाचे
निषेधनही भविष्यातील परिस्थितीवश विधान असते; म्हणून
$\tf{\lnot}(\Undef) = \Undef$. संधीमधील एक घटक अनिर्धारित
असेल आणि दुसरा सत्य असेल, तर ती संधीही अनिर्धारित असते.
कारण परिस्थितीवश घटक पुढे कसा ठरतो त्यानुसार संधी सत्यही
ठरू शकते आणि असत्यही ठरू शकते. म्हणून \[    \tf{\land}(\True, \Undef) =
\tf{\land}(\Undef, \True) =
\Undef.
\]
पण दुसरा घटक असत्य असेल, तर संधी सत्य ठरू शकत नाही;
म्हणून \[\tf{\land}(\False, \Undef) =
\tf{\land}(\Undef, \False) = \False.\]
दोन्ही निविष्टे $\True$ किंवा $\False$ ही निश्चित सत्यतामूल्ये
असतील, तर उरलेली फलनमूल्ये अभिजात तर्कशास्त्राप्रमाणेच असतात.

सशर्त विधानाच्या बाबतीत परिस्थिती थोडी अधिक गुंतागुंतीची आहे.
$q$ हे भविष्यातील परिस्थितीवश विधान आहे असे समजा. $p$ असत्य
असेल, तर $q$ पुढे कसेही ठरले तरी $p \lif q$ सत्य असेल; म्हणून
$\tf{\lif}(\False, \Undef) = \True$ असे ठेवावे. तसेच $p$
सत्य असेल, तर $q$ पुढे कसेही ठरले तरी $q \lif p$ सत्य असेल;
म्हणून $\tf{\lif}(\Undef, \True) = \True$. $p$ सत्य असेल,
तर $p \lif q$ सत्य किंवा असत्य ठरू शकते; म्हणून
$\tf{\lif}(\True, \Undef) = \Undef$. त्याचप्रमाणे $p$
असत्य असेल, तर $q \lif p$ सत्य किंवा असत्य ठरू शकते;
म्हणून $\tf{\lif}(\Undef, \False) = \Undef$. आता $p$
आणि $q$ दोन्ही भविष्यातील परिस्थितीवश विधाने असण्याचे
प्रकरण उरते. मूळ प्रेरणेनुसार या प्रकरणात खरे तर $\Undef$
हेच मूल्य द्यायला हवे. मात्र तसे केल्यास $\varphi \lif \varphi$
हे \emph{सर्वतः सत्य सूत्र राहणार नाही}. $\varphi \lor \lnot \varphi$
आणि $\lnot(\varphi \land \lnot \varphi)$ यांचा त्याग करायला
\L ukasiewicz यांना अडचण वाटली नाही; पण $\varphi \lif \varphi$
चा त्याग करणे त्यांना मान्य नव्हते. म्हणून त्यांनी
$\tf{\lif}(\Undef, \Undef) = \True$ असे ठरवले.

\begin{defn}\label{mvl:thr:luk:def:lukasiewicz}
त्रिमूल्य \L ukasiewicz तर्कशास्त्र पुढील मॅट्रिक्सने परिभाषित होते:
\begin{enumerate}
  \item मानक विधानीय भाषा $\Lang L_0$ चा केवळ
  $\lnot$, $\land$, $\lor$, $\lif$ हे संयोजक असलेला खंड.
  \item सत्यतामूल्यांचा संच $V = \{\True, \Undef, \False\}$.
  \item $\True$ हेच एकमेव निर्दिष्ट मूल्य आहे, म्हणजे $V^+ = \{\True\}$.
  \item सत्यता-फलने पुढील कोष्टकांनी दिली आहेत:
  \begin{center}
    \begin{tabular}{c|c}
      $\tf{\lnot}$ & \\
      \hline
      $\True$ & $\False$ \\
      $\Undef$ & $\Undef$ \\
      $\False$ & $\True$
    \end{tabular}
    \quad
    \begin{tabular}{c|ccc}
      $\tf{\land}[\LogLuk[3]]$ & $\True$ & $\Undef$ & $\False$ \\
      \hline
      $\True$ & $\True$ & $\Undef$ & $\False$ \\
      $\Undef$ & $\Undef$ & $\Undef$ & $\False$\\
      $\False$ & $\False$ & $\False$ & $\False$
    \end{tabular}
    \\[2ex]
    \begin{tabular}{c|ccc}
      $\tf{\lor}[\LogLuk[3]]$ & $\True$ & $\Undef$ & $\False$ \\
      \hline
      $\True$ & $\True$ & $\True$ & $\True$ \\
      $\Undef$ & $\True$ & $\Undef$ & $\Undef$ \\
      $\False$ & $\True$ & $\Undef$ & $\False$
    \end{tabular}
    \quad
    \begin{tabular}{c|ccc}
      $\tf{\lif}[\LogLuk[3]]$ & $\True$ & $\Undef$ & $\False$ \\
      \hline
      $\True$ & $\True$ & $\Undef$ & $\False$ \\
      $\Undef$ & $\True$ & $\True$ & $\Undef$  \\
      $\False$ & $\True$ & $\True$ & $\True$
    \end{tabular}
  \end{center}
\end{enumerate}
\end{defn}
सहज लक्षात येते की केवळ $\lnot$, $\land$ आणि $\lor$
असलेल्या कोणत्याही सूत्र~$\varphi$ मधील सर्व
विधानीय चले ना $\Undef$ नेमले, तर त्या
सूत्राचे सत्यतामूल्य~$\Undef$ असेल. म्हणून, उदाहरणार्थ,
अभिजात तर्कशास्त्रातील $p \lor \lnot p$ आणि $\lnot(p
\land \lnot p)$ ही सर्वतः सत्य सूत्रे $\LogLuk[3]$
मध्ये सर्वतः सत्य नाहीत; कारण $\pAssign v(p) = \Undef$
असताना $\pValue{v}(\varphi) = \Undef$.
संबंधित प्रत्येक विधानीय चलासाठी $\pAssign v(p) = \True$ किंवा
$\False$ असलेल्या सत्य-मूल्यांकने मध्ये $\pValue v(\varphi)$
हे त्या सूत्राच्या अभिजात सत्यतामूल्याशी जुळते.
\begin{prop}
  $\varphi$ मधील प्रत्येक $p$ साठी $\pAssign v(p) \in \{\True, \False\}$ असेल, तर
  $\pValue v(\varphi)[\LogLuk[3]] = \pValue v(\varphi)[\LogCL]$.
\end{prop}
\begin{prob}\label{mvl:thr:luk:prob:luk-iff} असे समजा की आपण
  $\pValue v(\varphi \liff \psi) = \pValue v((\varphi \lif \psi) \land (\psi \lif
  \varphi))$ अशी $\LogLuk[3]$ मध्ये व्याख्या करतो. मग $\liff$
  चे सत्यताकोष्टक कसे असेल?
\end{prob}
अभिजात तर्कशास्त्रातील अनेक सर्वतः सत्य सूत्रे \LogLuk[3]
मध्येही \emph{सर्वतः सत्य आहेत}; उदाहरणार्थ,
$\lnot p \lif (p \lif q)$. अभिजात तर्कशास्त्राप्रमाणेच
हे सत्यताकोष्टकाने पडताळता येते:
\begin{center}
\begin{tabular}{cc|cccccc}
  $p$ & $q$ & $\lnot$ & $p$ & $\lif$ & $\smash{(}p$ & $\lif$ & $q\smash{)}$ \\ \hline
  \True & \True & \False & \True & \True & \True & \True & \True \\
  \True & \Undef & \False & \True & \True & \True & \Undef & \Undef \\
  \True & \False & \False & \True & \True & \True & \False & \False \\
  \Undef & \True & \Undef & \Undef & \True & \Undef & \True & \True \\
  \Undef & \Undef & \Undef & \Undef & \True & \Undef & \True & \Undef \\
  \Undef & \False & \Undef & \Undef & \True & \Undef & \Undef & \False \\
  \False & \True & \True & \False & \True & \False & \True & \True \\
  \False & \Undef & \True & \False & \True & \False & \True & \Undef \\
  \False & \False & \True & \False & \True & \False & \True & \False \\
\end{tabular}
\end{center}
\begin{prob}
  पुढील सूत्रे \LogLuk[3] मध्ये सर्वतः सत्य आहेत हे दाखवा:
  \begin{enumerate}
    \item $p \lif (q \lif p)$
    \item\label{mvl:thr:luk:prob:luk-taut-2} $\lnot(p \land q) \liff (\lnot p \lor \lnot q)$
    \item\label{mvl:thr:luk:prob:luk-taut-3} $\lnot(p \lor q) \liff (\lnot p \land \lnot q)$
  \end{enumerate}
  (\ref{mvl:thr:luk:prob:luk-taut-2} आणि
  \ref{mvl:thr:luk:prob:luk-taut-3} मध्ये $\varphi \liff \psi$
  हा $(\varphi \lif \psi) \land (\psi \lif \varphi)$ चा संक्षेप माना किंवा
  \ref{mvl:thr:luk:prob:luk-iff} साठीचे तुमचे उत्तर वापरा.)
\end{prob}
\begin{prob}
  पुढील अभिजात सर्वतः सत्य सूत्रे \LogLuk[3] मध्ये सर्वतः सत्य
  नाहीत हे दाखवा:
  \begin{enumerate}
    \item $(\lnot p \land p) \lif q$
    \item $((p \lif q) \lif p) \lif p$
    \item $(p \lif (p \lif q)) \lif (p \lif q)$
  \end{enumerate}
\end{prob}
त्यामुळे अभिजात तर्कशास्त्रातील सर्व सर्वतः सत्य सूत्रे
$\LogLuk[3]$ मध्ये सर्वतः सत्य नसली, तरी प्रत्येक
सत्य-मूल्यांकन मध्ये त्यांना निदान $\True$ किंवा $\Undef$
यांपैकी एक मूल्य मिळेल, असे वाटू शकते. पण तसे नाही.
पुढील सूत्र हे प्रतिउदाहरण आहे:
\[
  \lnot(p \lif \lnot p) \lor \lnot(\lnot p \lif p)
\]
$p$ चे मूल्य~$\Undef$ असेल, तर या सूत्राचे मूल्य $\False$ येते.

\begin{prob}
  पुढीलपैकी कोणते संबंध \L ukasiewicz तर्कशास्त्रात खरे ठरतात?
  प्रत्येकासाठी सत्यताकोष्टक द्या.
  \begin{enumerate}
    \item $p, p \lif q \Entails q$
    \item $\lnot\lnot p \Entails p$
    \item $p \land q \Entails p$
    \item $p \Entails p \land p$
    \item $p \Entails p \lor q$
  \end{enumerate}
\end{prob}

\L ukasiewicz यांना आपल्या त्रिमूल्य पद्धतीच्या आधारे
संभाव्यतेचे तर्कशास्त्र उभारायचे होते. त्यासाठी त्यांनी ``$\varphi$
संभाव्य आहे'' याकरिता एकस्थानी संयोजक $\Diamond \varphi$ आणि
``$\varphi$ अनिवार्य आहे'' याकरिता त्याला अनुरूप $\Box \varphi$ घेतले:
\begin{center}
  \begin{tabular}{c|c}
    $\tf{\Diamond}$ & \\
    \hline
    $\True$ & $\True$ \\
    $\Undef$ & $\True$ \\
    $\False$ & $\False$
  \end{tabular}
  \quad
  \begin{tabular}{c|c}
    $\tf{\Box}$ & \\
    \hline
    $\True$ & $\True$ \\
    $\Undef$ & $\False$ \\
    $\False$ & $\False$
  \end{tabular}
\end{center}
म्हणजे $p$ संभाव्य आहे, जर आणि फक्त जर ते असत्य असल्याचे
आधीच निश्चित झालेले नसेल; आणि $p$ अनिवार्य आहे, जर आणि
फक्त जर ते सत्य असल्याचे आधीच निश्चित झालेले असेल.

\begin{prob}
  $\Box$ आणि~$\Diamond$ यांची सत्यताकोष्टके जोडून विस्तारलेल्या
  \LogLuk[3] मध्ये $\Box p \liff \lnot\Diamond \lnot p$ आणि
  $\Diamond p \liff \lnot \Box \lnot p$ ही सर्वतः सत्य सूत्रे
  आहेत हे दाखवा.
\end{prob}

पण या प्रस्तावित मोडल तर्कशास्त्राच्या मर्यादा लवकरच स्पष्ट
झाल्या: पुढे काहीही घडले तरी $p \land \lnot p$ कधीच सत्य
ठरू शकत नाही. म्हणून सध्या त्याचे मूल्य निश्चित नसले
(आणि त्यामुळे ते अनिर्धारित असले), तरी ते अशक्य मानायला
हवे; म्हणजे $\lnot \Diamond(p \land \lnot p)$ हे सर्वतः सत्य
सूत्र असायला हवे. तथापि, $\pAssign v(p) = \Undef$ असेल,
तर $\pValue v(\lnot \Diamond(p \land \lnot p)) = \False$.
कारण या वेळी $p \land \lnot p$ चे मूल्य $\Undef$,
त्याच्या $\Diamond$ चे मूल्य $\True$, आणि शेवटच्या निषेधनाचे
मूल्य $\False$ असते.
संभाव्यता आणि अनिवार्यता यांसारखे मोडल भेद व्यक्त करण्यासाठी
दोन सत्यतामूल्ये पुरेशी नाहीत, हा \L ukasiewicz यांचा निष्कर्ष
बरोबर होता; पण तिसरे सत्यतामूल्य जोडणेही पुरेसे नाही.













\section{क्लिनी तर्कशास्त्रे}\label{mvl:thr:skl:sec}

Stephen Kleene यांनी संगणन प्रक्रियांच्या परिणामांच्या दृष्टीने
सत्यतामूल्यांचा विचार करणाऱ्या तर्कशास्त्रातून प्रेरित होऊन दोन
त्रिमूल्य तर्कशास्त्रे मांडली: एखादी प्रक्रिया $\True$ किंवा
$\False$ हा परिणाम देऊ शकते; पण ती समाप्तही होणार नाही असे
घडू शकते. त्या वेळी संबंधित सत्यतामूल्य अपरिभाषित असते;
ते सत्यतामूल्य~$\Undef$ ने दर्शवले जाते.

विधान~$\varphi$ चे निषेधन संगणित करण्यासाठी आधी~$\varphi$ चे मूल्य
काढायचे आणि मग त्याच्या विरुद्ध मूल्याचा परिणाम द्यायचा.
$\varphi$ चे संगणन समाप्त होत नसेल, तर ही संपूर्ण प्रक्रियाही
समाप्त होत नाही; म्हणून $\Undef$ चे निषेधन $\Undef$ आहे.

संधी $\varphi \land \psi$ संगणित करण्याचे दोन मार्ग आहेत: आधी~$\varphi$
आणि नंतर $\psi$ संगणित करता येईल. दोन्हींचा परिणाम~$\True$
असेल तर संधीचा परिणाम $\True$, अन्यथा $\False$ असेल.
दोन्हींपैकी कोणतेही संगणन थांबले नाही, तर संपूर्ण प्रक्रियाही
थांबत नाही. त्यामुळे या पद्धतीत संधीचा एक घटक अपरिभाषित
असेल, तर संधीही अपरिभाषित असते. विसंधीबाबतही हेच लागू होते.

पण $\varphi$ आणि $\psi$ यांचे मूल्यांकन समांतरपणे करता आले, तर
अधिक चांगले करता येते. दोन प्रक्रियांपैकी एक थांबून $\False$
हा परिणाम देताच आपण थांबू शकतो, कारण संधी असत्य असलीच
पाहिजे. म्हणून या पद्धतीत एक घटक असत्य असेल, तर दुसरा
अपरिभाषित असला तरी संधी असत्य असते. त्याचप्रमाणे विसंधीचे
समांतर संगणन करताना दोन घटकांपैकी एक सत्य परिणाम देताच
थांबता येते; विसंधी सत्य असलीच पाहिजे. त्यामुळे संयुक्त
विधानाचा एक घटक अपरिभाषित असला तरी त्याचा परिणाम कळू
शकतो. या अर्थानुसार $\Undef$ ला ``अपरिभाषित'' याऐवजी
``अज्ञात'' असे वाचता येईल.

या दोन अर्थांमधून क्लिनी यांची प्रबळ आणि दुर्बल तर्कशास्त्रे
मिळतात. सशर्त विधानाची व्याख्या $\lnot \varphi \lor \psi$ शी
समतुल्य अशी केली जाते.

\begin{defn}
\emph{प्रबळ क्लिनी तर्कशास्त्र}~$\LogKs$ पुढील मॅट्रिक्सने
परिभाषित होते:
\begin{enumerate}
  \item मानक विधानीय भाषा $\Lang L_0$ चा केवळ
  $\lnot$, $\land$, $\lor$, $\lif$ हे संयोजक असलेला खंड.
  \item सत्यतामूल्यांचा संच $V = \{\True, \Undef, \False\}$.
  \item $\True$ हेच एकमेव निर्दिष्ट मूल्य आहे, म्हणजे $V^+ = \{\True\}$.
  \item सत्यता-फलने पुढील कोष्टकांनी दिली आहेत:
  \begin{center}
    \begin{tabular}{c|c}
      $\tf{\lnot}$ & \\
      \hline
      $\True$ & $\False$ \\
      $\Undef$ & $\Undef$ \\
      $\False$ & $\True$
    \end{tabular}
    \quad
    \begin{tabular}{c|ccc}
      $\tf{\land}[\LogKs]$ & $\True$ & $\Undef$ & $\False$ \\
      \hline
      $\True$ & $\True$ & $\Undef$ & $\False$ \\
      $\Undef$ & $\Undef$ & $\Undef$ & $\False$\\
      $\False$ & $\False$ & $\False$ & $\False$
    \end{tabular}
    \\[2ex]
    \begin{tabular}{c|ccc}
      $\tf{\lor}[\LogKs]$ & $\True$ & $\Undef$ & $\False$ \\
      \hline
      $\True$ & $\True$ & $\True$ & $\True$ \\
      $\Undef$ & $\True$ & $\Undef$ & $\Undef$ \\
      $\False$ & $\True$ & $\Undef$ & $\False$
    \end{tabular}
    \quad
    \begin{tabular}{c|ccc}
      $\tf{\lif}[\LogKs]$ & $\True$ & $\Undef$ & $\False$ \\
      \hline
      $\True$ & $\True$ & $\Undef$ & $\False$ \\
      $\Undef$ & $\True$ & $\Undef$ & $\Undef$  \\
      $\False$ & $\True$ & $\True$ & $\True$
    \end{tabular}
  \end{center}
\end{enumerate}
\end{defn}
\begin{defn}
\emph{दुर्बल क्लिनी तर्कशास्त्र}~$\LogKw$ पुढील मॅट्रिक्सने
परिभाषित होते:
\begin{enumerate}
  \item मानक विधानीय भाषा $\Lang L_0$ चा केवळ
  $\lnot$, $\land$, $\lor$, $\lif$ हे संयोजक असलेला खंड.
  \item सत्यतामूल्यांचा संच $V = \{\True, \Undef, \False\}$.
  \item $\True$ हेच एकमेव निर्दिष्ट मूल्य आहे, म्हणजे $V^+ = \{\True\}$.
  \item सत्यता-फलने पुढील कोष्टकांनी दिली आहेत:
  \begin{center}
    \begin{tabular}{c|c}
      $\tf{\lnot}$ & \\
      \hline
      $\True$ & $\False$ \\
      $\Undef$ & $\Undef$ \\
      $\False$ & $\True$
    \end{tabular}
    \quad
    \begin{tabular}{c|ccc}
      $\tf{\land}[\LogKw]$ & $\True$ & $\Undef$ & $\False$ \\
      \hline
      $\True$ & $\True$ & $\Undef$ & $\False$ \\
      $\Undef$ & $\Undef$ & $\Undef$ & $\Undef$\\
      $\False$ & $\False$ & $\Undef$ & $\False$
    \end{tabular}
    \\[2ex]
    \begin{tabular}{c|ccc}
      $\tf{\lor}[\LogKw]$ & $\True$ & $\Undef$ & $\False$ \\
      \hline
      $\True$ & $\True$ & $\Undef$ & $\True$ \\
      $\Undef$ & $\Undef$ & $\Undef$ & $\Undef$ \\
      $\False$ & $\True$ & $\Undef$ & $\False$
    \end{tabular}
    \quad
    \begin{tabular}{c|ccc}
      $\tf{\lif}[\LogKw]$ & $\True$ & $\Undef$ & $\False$ \\
      \hline
      $\True$ & $\True$ & $\Undef$ & $\False$ \\
      $\Undef$ & $\Undef$ & $\Undef$ & $\Undef$  \\
      $\False$ & $\True$ & $\Undef$ & $\True$
    \end{tabular}
  \end{center}
\end{enumerate}
\end{defn}
\begin{prop}
  $\LogKs$ आणि $\LogKw$ यांत कोणतेही सर्वतः सत्य सूत्र नाही.
\end{prop}
\begin{proof}
  सर्व विधानीय चले~$p$ साठी
  $\pAssign v(p) = \Undef$ असेल, तर कोणत्याही
  सूत्र~$\varphi$ चे सत्यतामूल्य~$\pValue v(\varphi) =
  \Undef$ असेल; कारण दोन्ही तर्कशास्त्रांत
  \[
    \tf{\lnot}(\Undef) = \tf{\lor}(\Undef, \Undef) = \tf{\land}(\Undef,
  \Undef) = \tf{\lif}(\Undef, \Undef) = \Undef
  \]
  असे आहे. $\LogKs$ आणि $\LogKw$ या दोन्हींमध्ये
  $\Undef \notin V^+$ असल्यामुळे या सत्य-मूल्यांकन मध्ये
  $\varphi$ चे सत्यतामूल्य निर्दिष्ट नसेल.
\end{proof}

दुर्बल आणि प्रबळ क्लिनी तर्कशास्त्रांत सर्वतः सत्य सूत्रे
नसली, तरी त्यांचे तार्किक निष्पन्नता संबंध अतुच्छ आहेत.

\begin{prob}
  पुढीलपैकी कोणते संबंध (a) प्रबळ आणि (b) दुर्बल क्लिनी
  तर्कशास्त्रात खरे ठरतात? प्रत्येकासाठी सत्यताकोष्टक द्या.
  \begin{enumerate}
    \item $p, p \lif q \Entails q$
    \item $p \lor q, \lnot p \Entails q$
    \item $p \land q \Entails p$
    \item $p \Entails p \land p$
    \item $p \Entails p \lor q$
  \end{enumerate}
\end{prob}

Dmitry Bochvar यांनी $\Undef$ चा अर्थ ``अर्थहीन'' असा घेतला
आणि विरोधाभास निर्माण करणाऱ्या विधानांना~$\Undef$ हे मूल्य
देऊन खोटारड्याच्या विरोधाभासासारखे विरोधाभास सोडवण्याचा
प्रयत्न केला. त्यांनी मूलतः दुर्बल क्लिनी तर्कशास्त्राला
अतिरिक्त संयोजके जोडून तयार होणारे तर्कशास्त्र मांडले.
त्यांतील दोन संयोजके ``बाह्य निषेधन'' आणि ``अपरिभाषित
आहे'' हा परिकर्मी अशी आहेत:
\begin{center}
  \begin{tabular}{c|c}
    $\tf{\sim}$ & \\
    \hline
    $\True$ & $\False$ \\
    $\Undef$ & $\True$ \\
    $\False$ & $\True$
  \end{tabular}
\quad
  \begin{tabular}{c|c}
    $\tf{+}$ & \\
    \hline
    $\True$ & $\False$ \\
    $\Undef$ & $\True$ \\
    $\False$ & $\False$
  \end{tabular}
\end{center}

\begin{prob}
  Bochvar यांच्या तर्कशास्त्रात $\lnot$ आणि~$+$ यांच्या
  साहाय्याने $\sim$ ची व्याख्या करता येते का? म्हणजे, फक्त
  विधानीय चल~$p$ असलेले, $\sim$ नसलेले
  आणि नेहमी~$\mathord{\sim}p$ इतकेच सत्यतामूल्य घेणारे
  सूत्र शोधा. तुमचे उत्तर योग्य आहे हे
  सत्यताकोष्टकाने दाखवा.
\end{prob}













\section{G\"odel तर्कशास्त्रे}\label{mvl:thr:god:sec}

Kurt G\"odel यांनी $n$-मूल्य तर्कशास्त्रांची एक क्रमिका
मांडली. त्यांतील प्रत्येकात अंतःप्रज्ञावादी तर्कशास्त्रात
वैध असलेली सर्व सूत्रे येतात आणि प्रत्येक अभिजात
तर्कशास्त्राचे उपतर्कशास्त्र आहे. त्यांपैकी पहिले
लक्षवेधी तर्कशास्त्र पुढीलप्रमाणे आहे:

\begin{defn}\label{mvl:thr:god:defn:goedel}
\emph{$3$-मूल्य G\"odel तर्कशास्त्र}~$\LogGod$ पुढील
मॅट्रिक्सने परिभाषित होते:
\begin{enumerate}
  \item $\lfalse$, $\lnot$, $\land$, $\lor$, $\lif$ असलेली मानक विधानीय भाषा $\Lang L_0$.
  \item सत्यतामूल्यांचा संच $V = \{\True, \Undef, \False\}$.
  \item $\True$ हेच एकमेव निर्दिष्ट मूल्य आहे, म्हणजे $V^+ = \{\True\}$.
  \item $\lfalse$ साठी $\tf{\lfalse} = \False$. उरलेल्या
  संयोजकांची सत्यता-फलने पुढील कोष्टकांनी दिली आहेत:
  \begin{center}
    \begin{tabular}{c|c}
      $\tf{\lnot}[\LogGod]$ & \\
      \hline
      $\True$ & $\False$ \\
      $\Undef$ & $\False$ \\
      $\False$ & $\True$
    \end{tabular}
    \quad
    \begin{tabular}{c|ccc}
      $\tf{\land}[\LogGod]$ & $\True$ & $\Undef$ & $\False$ \\
      \hline
      $\True$ & $\True$ & $\Undef$ & $\False$ \\
      $\Undef$ & $\Undef$ & $\Undef$ & $\False$\\
      $\False$ & $\False$ & $\False$ & $\False$
    \end{tabular}
    \\[2ex]
    \begin{tabular}{c|ccc}
      $\tf{\lor}[\LogGod]$ & $\True$ & $\Undef$ & $\False$ \\
      \hline
      $\True$ & $\True$ & $\True$ & $\True$ \\
      $\Undef$ & $\True$ & $\Undef$ & $\Undef$ \\
      $\False$ & $\True$ & $\Undef$ & $\False$
    \end{tabular}
    \quad
    \begin{tabular}{c|ccc}
      $\tf{\lif}[\LogGod]$ & $\True$ & $\Undef$ & $\False$ \\
      \hline
      $\True$ & $\True$ & $\Undef$ & $\False$ \\
      $\Undef$ & $\True$ & $\True$ & $\False$  \\
      $\False$ & $\True$ & $\True$ & $\True$
    \end{tabular}
  \end{center}
\end{enumerate}
\end{defn}
$\land$ आणि~$\lor$ यांची सत्यताकोष्टके \L ukasiewicz
आणि प्रबळ क्लिनी तर्कशास्त्रांप्रमाणेच आहेत, हे लक्षात येईल;
पण $\lnot$ आणि~$\lif$ यांच्या कोष्टकांत फरक आहे.
G\"odel तर्कशास्त्रात $\tf{\lnot}(\Undef) = \False$.
\L ukasiewicz आणि क्लिनी तर्कशास्त्रांच्या विपरीत,
$\tf{\lif}(\Undef, \False) = \False$; तर क्लिनी
तर्कशास्त्राच्या विपरीत (पण \L ukasiewicz तर्कशास्त्राप्रमाणे)
$\tf{\lif}(\Undef, \Undef) = \True$.
वर सूचित केलेल्या अंतःप्रज्ञावादी तर्कशास्त्राशी असलेल्या
संबंधावरून अपेक्षित आहे त्याप्रमाणे $\LogGod[3]$ हे
अंतःप्रज्ञावादी तर्कशास्त्राच्या जवळचे आहे. त्या तर्कशास्त्रात
वैध असलेली सर्व सूत्रे $\LogGod[3]$ मध्ये सर्वतः सत्य
असतात; आणि अंतःप्रज्ञावादी तर्कशास्त्रात वैध नसलेली अनेक
अभिजात सर्वतः सत्य सूत्रे $\LogGod[3]$ मध्येही सर्वतः
सत्य नसतात. उदाहरणार्थ, पुढील सूत्रे सर्वतः सत्य नाहीत:
\begin{align*}
  & p \lor \lnot p && (p \lif q) \lif (\lnot p \lor q) \\
  & \lnot\lnot p \lif p && \lnot(\lnot p \land \lnot q) \lif (p \lor q) \\
  & ((p \lif q) \lif p) \lif p && \lnot(p \lif q) \lif (p \land \lnot q)
\end{align*}
मात्र $\LogGod[3]$ मधील प्रत्येक सर्वतः सत्य सूत्र
अंतःप्रज्ञावादी तर्कशास्त्रातही वैध असेलच असे नाही;
उदाहरणार्थ, $\lnot\lnot p \lor \lnot p$ किंवा $(p
\lif q) \lor (q \lif p)$.
\begin{prob}
  पुढील सूत्रे $\LogGod[3]$ मध्ये सर्वतः सत्य आहेत हे
  सत्यताकोष्टकांनी दाखवा:
  \begin{align*}
    & \lnot\lnot p \lor \lnot p\\
    & (p \lif q) \lor (q \lif p) \\
    & \lnot(p \land q) \lif (\lnot p \lor \lnot q) \\
    & (p \lif q) \lor (q \lif r) \lor (r \lif s)
  \end{align*}
\end{prob}
\begin{prob}
  पुढील सूत्रे $\LogGod[3]$ मध्ये सर्वतः सत्य नाहीत हे
  सत्यताकोष्टकांनी दाखवा:
  \begin{align*}
    & (p \lif q) \lif (\lnot p \lor q) \\
    & \lnot(\lnot p \land \lnot q) \lif (p \lor q) \\
    & ((p \lif q) \lif p) \lif p \\
    & \lnot(p \lif q) \lif (p \land \lnot q)
  \end{align*}
\end{prob}
\begin{prob}
  पुढीलपैकी कोणते संबंध G\"odel तर्कशास्त्रात खरे ठरतात?
  प्रत्येकासाठी सत्यताकोष्टक द्या.
  \begin{enumerate}
    \item $p, p \lif q \Entails q$
    \item $p \lor q, \lnot p \Entails q$
    \item $p \land q \Entails p$
    \item $p \Entails p \land p$
    \item $p \Entails p \lor q$
  \end{enumerate}
\end{prob}
\section[इतर सत्यतामूल्येही निर्दिष्ट करणे]{$\True$ व्यतिरिक्तही मूल्य निर्दिष्ट करणे}\label{mvl:thr:mul:sec}
आतापर्यंत पाहिलेल्या सर्व तर्कशास्त्रांत निर्दिष्ट
सत्यतामूल्यांचा संच $V^+ = \{\True\}$ होता; म्हणजे एखाद्या
विधान निर्दिष्ट मानले जाई, जर आणि फक्त जर त्याचे
सत्यतामूल्य~$\True$ असेल.
पण एखादे विधान फक्त~$\False$ नसले, तरी ते सत्य मानता येईल.
अशा वेळी, उदाहरणार्थ, मॅट्रिक्समध्ये
$V^+ = \{\True, \Undef\}$ असे ठरवून तर्कशास्त्र मिळेल.
\begin{defn}
\emph{विरोधापत्तीचे तर्कशास्त्र}~$\LogLP$ पुढील मॅट्रिक्सने
परिभाषित होते:
\begin{enumerate}
  \item मानक विधानीय भाषा $\Lang L_0$ चा केवळ
  $\lnot$, $\land$, $\lor$, $\lif$ हे संयोजक असलेला खंड.
  \item सत्यतामूल्यांचा संच $V = \{\True, \Undef, \False\}$.
  \item $\True$ आणि $\Undef$ निर्दिष्ट आहेत, म्हणजे $V^+ = \{\True, \Undef\}$.
  \item सत्यता-फलने प्रबळ क्लिनी तर्कशास्त्रातील फलनांसारखीच आहेत.
\end{enumerate}
\end{defn}
\begin{defn}
Halld\'en यांचे \emph{निरर्थकतेचे तर्कशास्त्र}~$\LogHal$
पुढील मॅट्रिक्सने परिभाषित होते:
\begin{enumerate}
  \item मानक विधानीय भाषा $\Lang L_0$ च्या
  $\lnot$, $\land$, $\lor$, $\lif$ या संयोजकांच्या खंडाचा
  $1$-स्थानी संयोजक~$+$ जोडून केलेला विस्तार.
  \item सत्यतामूल्यांचा संच $V = \{\True, \Undef, \False\}$.
  \item $\True$ आणि $\Undef$ निर्दिष्ट आहेत, म्हणजे $V^+ = \{\True, \Undef\}$.
  \item सत्यता-फलने दुर्बल क्लिनी तर्कशास्त्रातील फलनांसारखीच
  आहेत; त्याखेरीज ``अर्थहीन आहे'' हा परिकर्मी आहे:
  \begin{center}
    \begin{tabular}{c|c}
    $\tf{+}$ & \\
    \hline
    $\True$ & $\False$ \\
    $\Undef$ & $\True$ \\
    $\False$ & $\False$
    \end{tabular}
  \end{center}
\end{enumerate}
\end{defn}
यांची सत्यताकोष्टके क्लिनी तर्कशास्त्रांशी जुळतात; पण
त्यांच्या विपरीत यांत सर्वतः सत्य सूत्रे \emph{आहेत}.
\begin{prop}\label{mvl:thr:mul:prop:LP-taut-CL} $\LogLP$ मधील सर्वतः सत्य
  सूत्रे आणि अभिजात विधानीय तर्कशास्त्रातील सर्वतः सत्य सूत्रे
  एकच आहेत.
\end{prop}
\begin{proof}
  \ref{mvl:syn:sub:prop:mvl-cl} नुसार $\Entails[\LogLP] \varphi$
  असेल, तर $\Entails[\LogCL] \varphi$. उलट दिशा दाखवण्यासाठी
  आपण सिद्ध करू की $\pValue v(\varphi)[\LogKs] = \False$ असेल असे
  सत्य-मूल्यांकन~$\pAssign v\colon \PVar \to \{\False, \True,
  \Undef\}$ असल्यास $\pValue {v'}(\varphi)[\LogCL] = \False$
  असे सत्य-मूल्यांकन~$\pAssign {v'}\colon \PVar \to \{\False, \True\}$
  असते. यावरून $\LogLP$ साठी अपेक्षित निष्कर्ष मिळतो. कारण
  $\LogKs$ आणि $\LogLP$ यांची वैशिष्ट्यपूर्ण सत्यता-फलने
  सारखीच आहेत आणि $\LogLP$ मध्ये फक्त $\False$ हेच मूल्य
  निर्दिष्ट नाही ($\LogLP$ आणि~$\LogKs$ यांतील तोच एक फरक
  आहे). म्हणून $\Entails/[\LogLP] \varphi$ असेल, तर कोणत्यातरी
  सत्य-मूल्यांकन~$\pAssign v$ साठी $\pValue v(\varphi)[\LogLP] = \pValue
  v[\LogKs](\varphi) = \False$. सिद्ध करावयाच्या विधानानुसार
  $\pValue{v'}[\LogCL](\varphi) = \False$, म्हणजे $\Entails/[\LogCL] \varphi$.
  हे सिद्ध करण्यासाठी आधी $\pAssign {v'}$ ची व्याख्या अशी करू:
  \[
    \pAssign {v'}(p) =
    \begin{cases}
    \True & \text{जर } \pAssign {v}(p) \in \{\True, \Undef\}\\
    \False & \text{इतर वेळी}
  \end{cases}
  \]
  आता $\varphi$ वरील विगमनाने दाखवू की (a)~$\pValue v(\varphi)[\LogKs]
  = \False$ असेल, तर $\pValue {v'}(\varphi)[\LogCL] = \False$;
  आणि (b)~$\pValue v(\varphi)[\LogKs] = \True$ असेल, तर
  $\pValue {v'}(\varphi)[\LogCL] = \True$.
  \begin{enumerate}
    \item विगमनाचा आधार: $\varphi \ident p$. (a) किंवा (b) चा
    पूर्वपक्ष लागू असेल, तेव्हा मूळ मूल्य अनुक्रमे असत्य किंवा
    सत्य असते; त्या प्रकरणांतच
    \ref{mvl:syn:val:defn:pValue} नुसार $\pValue v(\varphi)[\LogKs]
    = \pAssign v(p) = \pValue {v'}(\varphi)[\LogCL]$ ही समानता
    मिळते, आणि त्यावरून (a) व~(b) दोन्ही सिद्ध होतात.

    विगमनाच्या पायरीसाठी पुढील प्रकरणे पाहू:
    \item $\varphi \ident \lnot \psi$.
    \begin{enumerate}
      \item $\pValue v(\lnot\psi)[\LogKs] = \False$ असे समजा.
      $\tf{\lnot}[\LogKs]$ च्या व्याख्येनुसार
      $\pValue v(\psi)[\LogKs]  = \True$. विगमन गृहितकातील
      प्रकरण~(b) नुसार $\pValue {v'}(\psi)[\LogCL]  = \True$;
      म्हणून $\pValue {v'}(\lnot \psi)[\LogCL] = \False$.
      \item $\pValue v(\lnot\psi)[\LogKs] = \True$ असे समजा.
      $\tf{\lnot}[\LogKs]$ च्या व्याख्येनुसार
      $\pValue v(\psi)[\LogKs]  = \False$. विगमन गृहितकातील
      प्रकरण~(a) नुसार $\pValue {v'}(\psi)[\LogCL]  = \False$;
      म्हणून $\pValue {v'}(\lnot \psi)[\LogCL] = \True$.
    \end{enumerate}

    \item $\varphi \ident (\psi \land \chi)$.
    \begin{enumerate}
      \item $\pValue v(\psi \land \chi)[\LogKs] = \False$ असे समजा.
      $\tf{\land}[\LogKs]$ च्या व्याख्येनुसार
      $\pValue v(\psi)[\LogKs]  = \False$ किंवा
      $\pValue v(\chi)[\LogKs]  = \False$. विगमन गृहितकातील
      प्रकरण~(a) नुसार $\pValue {v'}(\psi)[\LogCL]  = \False$
      किंवा $\pValue {v'}(\chi)[\LogCL]  = \False$;
      म्हणून $\pValue {v'}(\psi \land \chi)[\LogCL] = \False$.
      \item $\pValue v(\psi \land \chi)[\LogKs] = \True$ असे समजा.
      $\tf{\land}[\LogKs]$ च्या व्याख्येनुसार
      $\pValue v(\psi)[\LogKs]  = \True$ आणि
      $\pValue v(\chi)[\LogKs]  = \True$. विगमन गृहितकातील
      प्रकरण~(b) नुसार $\pValue {v'}(\psi)[\LogCL]  = \True$
      आणि $\pValue {v'}(\chi)[\LogCL]  = \True$;
      म्हणून $\pValue {v'}(\psi \land \chi)[\LogCL] = \True$.
    \end{enumerate}
  \end{enumerate}

  उरलेली दोन प्रकरणे याच प्रकारची आहेत आणि सरावासाठी
  सोडली आहेत. दुसरा मार्ग असा: वरील पुराव्याने फक्त $\lnot$
  आणि~$\land$ असलेल्या सर्व सूत्रे साठी निष्कर्ष
  सिद्ध केला आहे. आता $\LogKs$ आणि $\LogCL$ या दोन्हींमध्ये
  कोणत्याही $\pAssign v$ साठी $\pValue v(\psi \lor
  \chi) = \pValue v(\lnot(\lnot\psi \land \lnot \chi))$ आणि $\pValue v(\psi \lif
  \chi) = \pValue v(\lnot(\psi \land \lnot \chi))$ या समानता
  वापरता येतात.
\end{proof}

\begin{prob}
  \ref{mvl:thr:mul:prop:LP-taut-CL} चा पुरावा पूर्ण
  करा; म्हणजे $\varphi \ident (\psi \lor \chi)$ आणि
  $\varphi \ident (\psi \lif \chi)$ या प्रकरणांसाठी (a) व~(b)
  सिद्ध करा.
\end{prob}

\begin{prob}
प्रत्येक अभिजात सर्वतः सत्य सूत्र $\LogHal$ मध्येही
सर्वतः सत्य आहे हे सिद्ध करा.
\end{prob}

दोन्ही तर्कशास्त्रांच्या समान चार-संयोजकीय भाषाखंडात त्यांची
सर्वतः सत्य सूत्रे अभिजात तर्कशास्त्रासारखीच असली, तरी त्यांचे
तार्किक निष्पन्नता संबंध वेगळे आहेत. Halld\'en यांच्या अतिरिक्त
संयोजकाचा समावेश असलेल्या सूत्रांसाठी हे समानतेचे विधान नाही. उदाहरणार्थ,
$\LogLP$ हे \emph{व्याघातसहिष्णु} आहे: त्यात
$\lnot p, p \Entails/ q$. म्हणून
$\lnot \varphi, \varphi \Entails \psi$ हे स्फोटाचे तत्त्व सर्वसाधारणपणे
लागू होत नाही. ($\varphi$ आणि $\psi$ यांच्या काही प्रकरणांत
ते लागू होते; उदाहरणार्थ, $\psi$~हे सर्वतः सत्य सूत्र असेल तर.)

\begin{prob}
  पुढीलपैकी कोणते संबंध (a)~$\LogLP$ आणि (b)~$\LogHal$
  मध्ये खरे ठरतात? प्रत्येकासाठी सत्यताकोष्टक द्या.
  \begin{enumerate}
    \item $p, p \lif q \Entails q$
    \item $\lnot q, p \lif q \Entails \lnot p$
    \item $p \lor q, \lnot p \Entails q$
    \item $\lnot p, p \Entails q$
    \item $p \Entails p \lor q$
    \item $p \lif q, q\lif r \Entails p \lif r$
  \end{enumerate}
\end{prob}

$\LogLuk[3]$ मध्ये $\Undef$ हे मूल्यही निर्दिष्ट मानल्यास
काय होईल?
केवळ निर्दिष्ट मूल्यांचा हा बदल केल्याने पुढील R-Mingle
मॅट्रिक्स मिळत नाही; त्यासाठी सशर्त विधानाचे सत्यता-फलनही बदलते.

\begin{defn}
  \emph{त्रिमूल्य R-Mingle}~$\LogRM[3]$ चा असत्यता-स्थिरासह
  विस्तार पुढील मॅट्रिक्सने परिभाषित होतो:
  \begin{enumerate}
    \item $\lfalse$, $\lnot$, $\land$, $\lor$, $\lif$ असलेली मानक विधानीय भाषा $\Lang L_0$.
    \item सत्यतामूल्यांचा संच $V = \{\True, \Undef, \False\}$.
    \item $\True$ आणि $\Undef$ निर्दिष्ट आहेत, म्हणजे $V^+ = \{\True, \Undef\}$.
    \item $\tf{\lfalse}=\False$. $\lnot$, $\land$ आणि $\lor$
    यांची सत्यता-फलने \L ukasiewicz तर्कशास्त्र~$\LogLuk[3]$
    मधील फलनांसारखीच आहेत. सशर्त विधानाचे सत्यताकोष्टक असे आहे:
    \begin{center}
      \begin{tabular}{c|ccc}
        $\tf{\lif}[\LogRM[3]]$ & $\True$ & $\Undef$ & $\False$ \\
        \hline
        $\True$ & $\True$ & $\False$ & $\False$ \\
        $\Undef$ & $\True$ & $\Undef$ & $\False$ \\
        $\False$ & $\True$ & $\True$ & $\True$
      \end{tabular}
    \end{center}
  \end{enumerate}
  येथे $\Undef$ हे मधले मूल्य सत्य आणि असत्य दोन्ही असल्याचे
  दर्शवते. सशर्त विधानाचे कोष्टक \L ukasiewicz यांच्या कोष्टकापेक्षा
  वेगळे आहे.\footnote{Barteld Kooi आणि Allard Tamminga,
  ``Completeness via Correspondence for Extensions of the Logic of
  Paradox,'' (2012), प्रथम ऑनलाइन आवृत्ती पृ. 10, विभाग 4;
  \url{https://doi.org/10.1017/S1755020312000196}.}
\end{defn}

\begin{prob}
  पुढीलपैकी कोणते संबंध $\LogRM[3]$ मध्ये खरे ठरतात?
  \begin{enumerate}
    \item $p, p \lif q \Entails q$
    \item $p \lor q, \lnot p \Entails q$
    \item $\lnot p, p \Entails q$
    \item $p \Entails p \lor q$
  \end{enumerate}
\end{prob}

निर्दिष्ट मूल्ये बदलल्यामुळे वेगवेगळी सत्यताकोष्टके कधी कधी
एकच तर्कशास्त्र, म्हणजे एकच तार्किक निष्पन्नता संबंध, देऊ
शकतात. उदाहरणार्थ, G\"odel तर्कशास्त्रात~$\{\True\}$ ऐवजी
$V^+ = \{\True, \Undef\}$ घेतल्यास असे घडते.

\begin{prop}\label{mvl:thr:mul:prop:gl-udes}
  $V = \{\False, \Undef,\True\}$,
  $V^+=\{\True, \Undef\}$ आणि $3$-मूल्य G\"odel
  तर्कशास्त्राची सत्यता-फलने असलेला मॅट्रिक्स अभिजात
  तर्कशास्त्र परिभाषित करतो.
\end{prop}

\begin{proof}
  सराव.
\end{proof}

\begin{prob}
  G\"odel तर्कशास्त्राप्रमाणेच पण
  $V^+=\{\True,\Undef\}$ घेऊन परिभाषित केलेल्या
  तर्कशास्त्र~$\Log L$ साठी $\Gamma \Entails/[\Log L] \psi$
  असेल, तर $\Gamma \Entails/[\LogCL] \psi$ हे दाखवून
  \ref{mvl:thr:mul:prop:gl-udes} सिद्ध करा.
  \ref{mvl:thr:mul:prop:LP-taut-CL} मधील कल्पना
  वापरा; मात्र गुणधर्म (a) आणि~(b) सिद्ध करण्याऐवजी
  $\pValue v(\varphi)[\LogGod] = \False$ जर आणि फक्त जर
  $\pValue {v'}(\varphi)[\LogCL] = \False$ हे दाखवा (आणि म्हणून
  $\pValue v(\varphi)[\LogGod] \in \{\True,\Undef\}$ जर आणि
  फक्त जर $\pValue {v'}(\varphi)[\LogCL] = \True$).
  यावरून विधान का सिद्ध होते ते स्पष्ट करा.
\end{prob}



